% Generated draft: requires mathematical and source review.
\documentclass[11pt]{article}
\usepackage[margin=25mm]{geometry}
\usepackage{fontspec}
\setmainfont{Latin Modern Roman}
\usepackage{amsmath,amssymb,mathtools,unicode-math}
\setmathfont{Latin Modern Math}
\usepackage{enumitem,needspace,xurl,hyperref}
\setlength{\parindent}{0pt}
\setlength{\parskip}{5pt}
\setlength{\emergencystretch}{3em}
\newcommand{\N}{\mathbb{N}}
\newcommand{\Z}{\mathbb{Z}}
\newcommand{\Q}{\mathbb{Q}}
\newcommand{\R}{\mathbb{R}}
\newcommand{\C}{\mathbb{C}}
\newcommand{\F}{\mathbb{F}}
\newcommand{\A}{\mathbb{A}}
\newcommand{\PP}{\mathbb{P}}
\newcommand{\E}{\mathbb{E}}
\newcommand{\eps}{\varepsilon}
\newcommand{\dd}{\,\mathrm{d}}
\begin{document}
\section*{UnsolvedMath ID: 30005900}
\textbf{Source problem code:} catalogue entry
\subsection{Definitions and mathematical statement}
Write \(X=M_{\mathrm{Dol}}(n,d)\) for the Dolbeault moduli space in the supplied record, where \(n\) denotes rank, \(d\) denotes degree, and \(g\) denotes the genus of the underlying curve. The target concerns singular spaces, including non-coprime cases. The record does not specify the underlying curve, the allowed genera, the structure group, determinant conditions, stability conventions, or the precise construction of this moduli space. These choices remain unresolved; the notation alone does not determine a unique moduli problem.

The following gives a conditional formulation of the intended correspondence. For a specified moduli problem admitting the required nonabelian Hodge correspondence, let \(Y=M_{\mathrm{B}}(n,d)\) denote its corresponding Betti moduli space, and let
\[
\Phi^j:IH^j(X;\Q)\longrightarrow IH^j(Y;\Q)
\]
be the induced identification. Here \(IH^j\) denotes middle-perversity intersection cohomology. Let \(P_k\) denote the perverse Leray filtration associated with the Hitchin map on the Dolbeault side, with the indexing intended by the conjecture, and let \(W_\ell\) denote the mixed Hodge weight filtration on the Betti side. Rational coefficients are used here to display the usual filtration formulation; the supplied record does not state a coefficient convention or a normalization of \(P\).

\textbf{Question (conditional filtration formulation).} For every singular Dolbeault moduli space in the intended family, including arbitrary integer degrees \(d\) for which that family is defined, is
\[
\Phi^j\!\left(P_k IH^j(X;\Q)\right)
 =W_{2k}IH^j(Y;\Q)
 =W_{2k+1}IH^j(Y;\Q)
\qquad(j,k\in\Z)?
\]
This is the intended \(PI=WI\) question, subject to specifying the moduli problem, correspondence, coefficients, and perverse indexing above. The displayed equality compares filtrations on the two sides of nonabelian Hodge theory; it does not place a weight filtration on Dolbeault intersection cohomology without that identification.

\textbf{Additional target.} Does curious hard Lefschetz hold on the corresponding Betti intersection cohomology for arbitrary degree \(d\)? To make the content of this target explicit, for a specified pure-dimensional Betti space of complex dimension \(D\), the usual weight-filtration form asks for a distinguished ordinary cohomology class \(\alpha\in H^2(Y;\Q)\) whose action raises weights by four, and isomorphisms
\[
\alpha^r\smile-:
\operatorname{Gr}^{W}_{D-2r}IH^j(Y;\Q)
\xrightarrow{\ \sim\ }
\operatorname{Gr}^{W}_{D+2r}IH^{j+2r}(Y;\Q)
\qquad(j\in\Z,\ r\in\Z_{\geq0}),
\]
where \(\operatorname{Gr}^{W}_{\ell}=W_\ell/W_{\ell-1}\). This is a conditional description of the customary target: the record does not supply the distinguished class, its normalization, or the hypotheses needed to apply this formulation to its unspecified family. The action here is an ordinary-cohomology action on intersection cohomology; it does not require a canonical ring structure on intersection cohomology itself.

\textbf{Supplied qualifications.} The report attributes the proposal of \(PI=WI\) to de Cataldo and Maulik in its reference [5, Question 4.1.7]. Leaving aside smooth cases, it reports proofs for \((d,g)=(0,1)\) and \((d,g,n)=(0,2,2)\) in [8, Main Thm]. It points to [17, \S5] for partial rank-two results and to [17, Corollary 1.5] for numerical evidence for curious hard Lefschetz, and describes arbitrary-degree curious hard Lefschetz as unknown at the time of the report. The supplied passage explains that ordinary cohomology of singular spaces need not satisfy Poincar\'e duality or relative hard Lefschetz, whereas intersection cohomology restores the relevant symmetries. It also emphasizes that intersection cohomology lacks a canonical ring structure, so proofs using multiplication in ordinary cohomology in the smooth setting do not automatically extend. The opening fragment, ``case \(d=0\)'', and the phrase ``all these cases'' refer to preceding context absent from the record and cannot specify additional hypotheses here.
\subsection{Short English statement}
The problem asks whether the perverse filtration on intersection cohomology of singular Dolbeault moduli spaces matches the weight filtration on the corresponding Betti spaces, and whether curious hard Lefschetz holds for arbitrary degree. The precise moduli family, filtration normalization, and Lefschetz class require clarification. The supplied status is partially solved; this is a statement-only draft requiring human mathematical review.
\subsection{Sources}
\begin{itemize}
\item Supplied local UnsolvedMath registry record, permanent ID \(30005900\), exact source problem code \texttt{catalogue entry}, title: \emph{Perverse-Weight Correspondence for Singular Dolbeault Moduli Spaces}. The record assigns proposed year 2024 and status \texttt{partially\_solved}. Its \texttt{corrected\_verified} label and statement that the formulation was checked against the official report are supplied registry assertions, not verification performed for this draft.
\item \emph{Arbeitsgemeinschaft: Geometry and Representation Theory around the P=W Conjecture} (2024), Oberwolfach Reports, pp.~949--1004, DOI 10.4171/owr/2024/16. Supplied report link: \url{https://doi.org/10.4171/owr/2024/16}. The record attributes the intersection-cohomology conjecture to de Cataldo and Maulik through the report's [5, Question 4.1.7]; complete bibliographic details for that numbered reference are not supplied.
\item Felisetti and Mauri, supplied reference label: \emph{PI=WI for symplectic resolutions}. \url{https://doi.org/10.5802/jep.196}.
\item Felisetti, supplied reference label: \emph{P=W phenomena survey}. \url{https://doi.org/10.1007/s40574-024-00420-z}.
\item The report passage additionally cites [8, Main Thm], [17, \S5], and [17, Corollary 1.5]. Their full bibliographic records and their correspondence with the separately supplied literature links are not specified in the record.
\item Literature assessment checked at \texttt{2026-08-22}: the supplied curation reports that the intersection version is proved in special cases, including symplectic-resolution cases, while the general singular and non-coprime case remains open. This is supplied dated curation derived from the OWR material, not an independently verified literature assessment. The record expressly requests independent verification of time-sensitive status and cited claims.
\end{itemize}
\end{document}
